\documentclass{article}
\renewcommand{\labelenumii}{\alph{enumii}}
\usepackage{enumerate}
\usepackage{listings}

\usepackage{graphicx,comment,framed}
\usepackage{roundbox}
\usepackage{fancybox}
\usepackage{tikz}
\usepackage{color}
\usepackage[hidelinks]{hyperref}
\usepackage{framed}
\usepackage{amsthm,amssymb,amsmath}
%\usepackage[colorlinks,linkcolor=blue,citecolor=blue]{hyperref}
\definecolor{shadecolor}{cmyk}{0,0,0,0}
\usepackage{listings}
\usepackage{xepersian}
\usepackage[noend]{algpseudocode}


%--------------------- page settings ----------------------

\settextfont[Scale=1.1]{XB Niloofar}
\setdigitfont[Scale=1.1]{XB Niloofar}
\defpersianfont\sayeh[Scale=1.1]{XB Niloofar}
\addtolength{\textheight}{3.2cm}
\addtolength{\topmargin}{-22mm}
\addtolength{\textwidth}{3cm}
\addtolength{\oddsidemargin}{-1.5cm}


%------------------------ Environments ------------------------------------

\newtheorem{قضیه}{قضیه}
\newtheorem{لم}{لم}
\newtheorem{مشاهده}{مشاهده}
\newtheorem{تعریف}{تعریف}


%-------------------------- Notations ------------------------------------

\newcommand{\IR}{\ensuremath{\mathbb{R}}} 
\newcommand{\IZ}{\ensuremath{\mathbb{Z}}} 
\newcommand{\IN}{\ensuremath{\mathbb{N}}} 
\newcommand{\IS}{\ensuremath{\mathbb{S}}} 
\newcommand{\IC}{\ensuremath{\mathbb{C}}} 
\newcommand{\IB}{\ensuremath{\mathbb{B}}} 

\newcommand{\bR}{\mathbb{R}}
\newcommand{\cB}{\mathcal{B}}
\newcommand{\cO}{\mathcal{O}}
\newcommand{\cG}{\mathcal{G}}
\newcommand{\rM}{\mathrm{M}}
\newcommand{\rC}{\mathrm{C}}
\newcommand{\rV}{\mathrm{V}}

\newcommand{\lee}{\leqslant}
\newcommand{\gee}{\geqslant}
\newcommand{\ceil}[1]{{\left\lceil{#1}\right\rceil}}
\newcommand{\floor}[1]{{\left\lfloor{#1}\right\rfloor}}
\newcommand{\prob}[1]{{\mbox{\tt Pr}[#1]}}
\newcommand{\set}[1]{{\{ #1 \}}}
\newcommand{\seq}[1]{{\left< #1 \right>}}
\newcommand{\provided}{\,|\,}
\newcommand{\poly}{\mbox{\rm poly}}
\newcommand{\polylog}{\mbox{\rm \scriptsize polylog}\,}
\newcommand{\comb}[2] {\left(\!\!\begin{array}{c}{#1}\\{#2}\end{array}\!\!\right)}




\newcounter{probcnt}
\newcommand{\مسئله}[1]{\stepcounter{probcnt}\bigskip\bigskip{
 	\large \bf مسئله‌ی \arabic{probcnt}$\mbox{\bf{.}}$ \ #1} \bigskip}

\newcommand{\fqed}[1]{\leavevmode\unskip\nobreak\quad\hspace*{\fill}{\ensuremath{#1}}}

\newenvironment{اثبات}
	{\begin{trivlist}\item[\hskip\labelsep{\em اثبات.}]}
	{\fqed{\square}\end{trivlist}}

\newenvironment{حل}
	{\begin{trivlist}\item[\hskip\labelsep{\bf حل.}]}
	{\fqed{\blacktriangleright}\end{trivlist}}

\ifdefined\hidesols
	\newsavebox{\trashcan} % uncomment the following line to hide solutions
	\renewenvironment{حل}{\begin{latin}\begin{lrbox}{\trashcan}}{\end{lrbox}\end{latin}}
\fi


%------------------------- Header -----------------------------

%------------------------- Header -----------------------------

\newcommand{\سربرگ}[3]{
	\parindent=0em
	\begin{shaded}
		
		\rightline{ 
			\makebox[6em][c]{
				\includegraphics[height=1.6cm]{aut.png}
		}}
		\vspace{-.2em}
		{\scriptsize\bf دانشکده‌ی مهندسی کامپیوتر}
		\hfill {\small
			مدرس: دکتر امیرمزلقانی  \ 
		}\\[-5em]
		\leftline{\hfill\Large\bf 
			کاربرد های جبر خطی
		}\\[.7em]
		\leftline{\hfill\bf 
			نیم‌سال دوم ۹۶-۹۷
		}\\[1.7em]
		\hrule height .12em
		
		\normalsize
		\vspace{1mm} #1
		\hfill \small  #3
		\vspace{1mm} 
		\hrule height .1em
		
		\vspace{-0.5em} 
		\hfill {\sayeh\large #2} \hfill
	\end{shaded}
	\begin{large}

توجه:
\end{large}
\\
\\
\begin{itemize}
\item
این تمرین از مباحث مربوط به فصل5 ،6و7 طراحی شده است که شامل 15سوال اجباری و 4 سوال امتیازی است که نمره سوال های امتیازی فقط به نمرات تمرین شما کمک می کند.
\item 
در برخی سوالات چند قسمتی از شما خواسته شده است به تعدادی از قسمت ها به اختیار پاسخ دهید ،پاسخگویی به بیش از مقدار تعیین شده نمره اضافی ندارد ،از بقیه قسمت ها می توانید برای تمرین بیشتر استفاده کنید.
\item 
در پایان نیز تعدادی سوال برای تمرین بیشتر در نظر گرفته شده است ،که به حل آن ها نمره ای تعلق نمی گیرد صرفا برای تمرین توسط خود شما در نظر گرفته شده است. 
\item 
اگه سوالی داشتین از طریق
\begin{latin}
\href{mailto:aut.la2018@gmail.com}{\[aut.la2018@gmail.com\]}
\end{latin}
 حتما بپرسید.

\item پاسخ های تمرین را در قالب یک فایل به صورت الگوی زیر آپلود کنید.
\begin{latin}
9531000\_steve\_McManaman\_HW4.pdf
\end{latin}
\item \color{red}
مهلت تحویل جمعه  1تیر 1397 ساعت 23:54:59
\end{itemize}

}




%\settextfont{Yas}
\begin{document}
\سربرگ{ فصل دوم (جبر ماتریسی)}{}{}
\clearpage

\مسئله{}  


الف) اگر $u = \begin{bmatrix}3 \\ 0\end{bmatrix}$ و $v = \begin{bmatrix}1\\ 2 \end{bmatrix}$ باشند. مساحت متواضی‌الاضلاع ساخته شده توسط $u$ و $v$ و $u + v$ و $0$ را بدست بیاورید.
\\\\
دترمینان $\begin{bmatrix}u & v\end{bmatrix}$ را بدست آورده و با مساحت متواضی‌الاضلاع مقایسه کنید. 
\\\\
درایه اول بردار $v$ را با یک مقدار دلخواه جایگزین کنید و تغیرات را مشاهده و تحیلل کنید.
\\\\
ب) نشان دهید معادله یک خط در ${\mathbb R}^2$ که از دو نقطه مشخص
$(x_1, y_1)$ و $(x_2, y_2)$میگذرد را میتوان به صورت زیر نوشت:

\[det(
\begin{bmatrix}
    1       & x        & y \\
    1       & x_1      & y_1 \\
    1       & x_2      & y_2 \\
\end{bmatrix}
)=0\]


\مسئله{}

معکوس ماتریس‌های زیر را به روش گوس جردن بدست بیاورید و مراحل آن را نیز بنویسید.
\\
\[A=
\begin{bmatrix}
    7       & 2      & 1 \\
    0       & 3      & -1 \\
    -3      & 4      & -2 \\
\end{bmatrix}
B=
\begin{bmatrix}
    1       & -2     & 1 \\
    4       & -7     & 3 \\
    -2      & 6      & -4 \\
\end{bmatrix}
\]

\مسئله{} 


درستی یا نادرستی گزاره های زیر را مشخص کنید و درصورت نادرست بودن مثال نقض ارائه دهید و در صورت درست بودن آن را اثبات کنید:
\begin{enumerate}
	\item اگر معکوس $A^2$ برابر $B$  باشد معکوس $A$ برابر $AB$ است.
\\
\item اگر $B$ و $C$ ماتریس‌هایی $m \times n$ باشند و $D$ معکوس پذیر باشد  و داشته باشیم 
$(B - C)D = 0$ ، $B = C$ است.
\\
\item اگر $A = BCD$ و داشته باشیم که $A$ معکوس‌پذیر است. $D$ و $C$ و $B$ هر سه معکوس پذیرند.
\\
\item اگر ماتریس $B$ یک ماتریس $3 \times 3$ باشد، $A = B^4 + 3B^2 + 7B + 3I_{3}$ معکوس پذیر است.

\end{enumerate}


\مسئله{}
به روش تجزیه $LU Factorization$ ماتریس زیر را تجزیه کرده و
ماتریس‌های $L$ و $U$ را به دست آورید. 
\\
\[A=
\begin{bmatrix}
    1       & 1      & 1 \\
    4       & 3      & -1 \\
    3       & 5      & 3 \\
\end{bmatrix}
\]

	\clearpage
{\LARGE{\textbf{سوالات شبیه سازی}}}

\مسئله

ماتریس متراکم (ماتریسی که بیشتر درایه های آن -مثلا بیش از نیمی از آنها- غیر صفر باشد) و معکوس پذیر
$A$
با ابعاد
$n \times n $
را در نظر بگیرید، روش استاندارد حل دستگاه معادله خطی 
$ Ax = b $
به صورت زیر است:
\begin{enumerate}
	\item 
	تجزیه
LU
	  ماتریس
A
	   را بیاید:
	$ A = LU$.
	
	\item 
	اگر 
	$\hat{x} := Ux$
	سیستم
	$L\hat{x} = b$
	(که در آن L یک ماتریس پایین مثلثی است) را از طریق جایگزینی پیشرو 
 $(forward substitution)$
حل کنید.
\item 
	سیستم بالا مثلثی 
$Ux= \hat{x} $
را(که در آن U یک ماتریس بالا مثلثی است) از طریق جایگزینی عقب گرد 
$(back substitution)$
حل کنید.
\end{enumerate}



\begin{latin}
	\begin{enumerate}[(a)]
		\begin{RTLitems}
			\item  \rl{
				تابعی بنویسید که تجزیه  LU ماتریس A را پیدا کند. فرض کنید که می توان ماتریس A را بدون استفاده از عمل جا به جایی دو سطر 
				$ row-interchange$
				از بین اعمال سطری مقدماتی
				به ماتریس بالا مثلثی
				$U$
				تبدیل کرد.
				تابع شما باید ماتریس A  را به عنوان ورودی بگیرد و ماتریس پایین مثلثی
				L
				و ماتریس بالا مثلثی
				U
				را باز گرداند.}
				\begin{latin}
				\begin{lstlisting}[language=Matlab]
Function [L, U] = lu_factor(A)
[n , n1] = size(A);
if n ~= n1 
	error ("A must be square")
end
L = eye (n)
U = zeros (n)

...

return;
				\end{lstlisting}
			\end{latin}
\rl{
				در کد بالا شما باید قسمت ....  را تکمیل کنید. برای این منظور تنها مقادیر بالای قطر اصلی ماتریس
				$U$
				که مقدار اولیه صفر گرفته است و مقادیر پایین قطر اصلی ماتریس L که برابر ماتریس همانی است را آپدیت کنید.	}
			

			\item  \rl{
			تابع دیگری بنویسید که معادله 
			Ax = b
			را از طریق مراحل ۱و ۲و ۳ را که در بالا ذکر شده است، حل کند. 
			تابع شما باید به شکل زیر باشد:}
			\begin{latin}
			\begin{lstlisting}[language=Matlab]

function x =  linear_sys_solver(A,b)

% compute the LU factorization of A
% Solve Ly = b for y by forward substitution
% Solve Ux = y by back substitution

return;
			\end{lstlisting}
		\end{latin}
			\rl{
	می توانید از کد خود مروبط به سوال ۸ تمرین اول در این بخش استفاده کنید	
}
			\item  \rl{
		تابع 
		myinverse
		را برای محاسبه وارون ماتریس 
		$A$
		با سایز 
		$n \times n $
		بنویسید. توجه کنید که باید از تابع 
		lu\_factor
		و توابع
		$forward substitution$
		و 
		$backward substitution$
		که در بخش های قبل یا تمرین قبل نوشته اید استفاده کنید.
		فرض کنید که 
		$ X = A^{-1}$
		پس
		$AX = I_n$
		از این نکته که 		
		$ Ax_i= e_i$ 
		به ازای
		$ i = 1 , ... , n$
		که در آن 
		$x_i$
		ستون 
		$i$
		ام ماتریس 
		$X$
		و 
		$e_{i}$
		ستون
		$i$
		ام ماتریس 
		$I_n$
		استفاده کنید.
		(توجه کنید که شما تنها یک بار می تواند از تجزیه 
		LU
		ماتریس A را محاسبه کنید)	
		}
	\item  \rl{
ماتریس هیلبرت یک ماتریس مربعی است به گونه ای که 
$H_{i,j} = \frac {1}{1+i+j}$
ماتریس هیلبرت A از مرتبه ۵ و ۱۰ و۱۵و ۲۰ را بسازید و ماتریس وارون آنها 
$(A^{-1})$
 را از طریق توابع آماده ( مثلا در متلب از طریق تابع 
inv
)
به دست آورید.
سپس ماتریس وارون آن را از طریق تابع 
myinverse 
 که در بخش قبل نوشته اید به دست آورید (آن را
 $A'^{-1}$
 بنامید)
 مقادیر 
 $AA^{-1}$
 و 
 $AA'^{-1}$
 را به دست آورید و نتایج را مقایسه و تحلیل کنید. (راهنمایی:‌ به ویژگی های ماتریس های 
 ill-conditioned
 توجه کنید
 )
}

		\end{RTLitems}    
	\end{enumerate}
\end{latin}



\end{document}
